\documentclass[a4paper,12pt]{report}
\usepackage[spanish]{babel}
\usepackage[T1]{fontenc}
\usepackage[utf8]{inputenc}
\usepackage{lmodern}
\usepackage{float}
\usepackage[margin=2.5cm]{geometry}
% ===== Numeración de secciones =====
\renewcommand\thesection{\arabic{section}}
\makeatletter
\makeatother
\usepackage[hidelinks]{hyperref}
% ===== Imagenes =====
\usepackage{graphicx}
\DeclareGraphicsExtensions{.png}
\graphicspath{{images/}}
% ===== Tablas =====
\usepackage{array}
\usepackage{tabularx}
\newcolumntype{Y}{>{\raggedright\arraybackslash}X}
% ===== Estilos para el recuadro bash para API =====
\usepackage{xcolor}
\usepackage{listings}
\definecolor{codebg}{RGB}{245,245,245}
\lstdefinestyle{bash}{
  language=bash,
  basicstyle=\ttfamily\small,
  backgroundcolor=\color{codebg},
  frame=single,
  breaklines=true,
  showstringspaces=false
}

\title{DOCUMENTEACIÓN\\\LARGE TRABAJO — CAFETERÍA MULTI-HILOS}
\author{Eder Martínez Castro}
\date{\today}

\begin{document}
\maketitle
\tableofcontents
\clearpage

\section{Análisis y Especificación de Requisitos}
\subsection{Descripción del problema}
\noindent Se requiere un \textbf{programa en Java} que \textbf{simule una cafetería} utilizando \textbf{hilos}, los \textbf{clientes} llegaran a la cafetería y serán atendidos por los \textbf{camareros} y los cafés los hara el \textbf{barista}.
\begin{itemize}
    \item Cada \textbf{cliente} tendrá \texttt{nombre} y \texttt{tiempoEspera}.
    \item Los \textbf{camareros} reciben los cafés del barista y lo entregan al cliente.
    \item El \textbf{barista} prapará los cafés.
    \item Si el timepo de espera del cliente se acaba, \textbf{se marchará}.
    \item El sistema terminará cuando \textbf{todos los clientes han recibido su cafe o se fueron}.
\end{itemize}

\subsection{Requisitos funcionales}
\noindent A continuación, se listan los requisitos funcionales del sistema:
\begin{table}[H]
    \centering
    \begin{tabularx}{\linewidth}{|l|Y|}
      \hline
      \textbf{ID} & \textbf{Descripción} \\
      \hline
      RF-01 & \textbf{Crear clientes}: \texttt{nombre:String} y \texttt{tiempoEspera:int (ms)}. \\
      \hline
      RF-02 & \textbf{Crear camareros}: instanciar $n$ camareros que serán los consumidores, cada uno se ejecutará en un hilo independiente. \\
      \hline
      RF-03 & \textbf{Crear barista}: instanciar un barista que será nustro productor. \\
      \hline
      RF-04 & \textbf{Preparar café}: el barista simulara con \texttt{Threads.sleep} el tiempo que tarda en prepararlo. \\
      \hline
      RF-05 & \textbf{Recoger café}: el camarero recogerá el café cuando el barista ya lo haya hecho. \\
      \hline
      RF-06 & \textbf{Abandono por espera}: si el cliente no recibe su café y se le acabo su \texttt{tiempoEspera}, se marchará. \\
      \hline
      RF-07 & \textbf{Entrega del café}: cuando el camarero le entregará el café al cliente si sigue esperando. \\
      \hline
    \end{tabularx}
\end{table}

\subsection{Requisitos no funcionales}
\noindent A continuación, se listan los requisitos no funcionales del sistema:
\begin{table}[H]
  \centering
  \begin{tabularx}{\linewidth}{|l|Y|}
    \hline
    \textbf{ID} & \textbf{Descripción} \\
    \hline
    RNF-01 & \textbf{Lenguaje}: El programa se realizará en \texttt{Java} y \texttt{JavaFX}. \\
    \hline
    RNF-02 & \textbf{Método a emplear}: Se utilizará el método \texttt{Productor-Consumidor}. \\
    \hline
    RNF-03 & \textbf{Observabilidad}: Los mensjes mostrados deben ser claros y simples. \\
    \hline
    RNF-04 & \textbf{Cierre ordenado}: todos los hilos terminaran de forma correcta  antes de que termine el programa. \\
    \hline
  \end{tabularx}
\end{table}

\subsection{Casos de uso}
\noindent \textbf{Descripción de casos de uso principales:}
\begin{table}[H]
  \centering
  \begin{tabularx}{\linewidth}{|l|Y|}
    \hline
    \textbf{Caso de uso} & \textbf{Descripción} \\
    \hline
    CU-01 Llegada y pedido & El \textbf{cliente} llega, y lo atiende el camarero y esperar su café hasta que se lo de o su \texttt{tiempoEspera} acabe. \\
    \hline
    CU-02 Preparación & El \textbf{barista} empieza a hacer cafés y notifica al camarero cuando ya esté listo. \\
    \hline
    CU-03 Coger el café cuando este listo & El \textbf{camarero} cuando le avisa el barista que el café esta listo lo recoge para entregarlo. \\
    \hline
    CU-04 Entrega & Si el cliente no se le termino el tiempo de espera, el camarero le \textbf{entrega el café} y se muestra el \textbf{tiempo total} de preparación de ese café por parte del camarero. \\
    \hline
    CU-05 Abandono & Si el tiempo de espera termina antes de la entrega, el cliente \textbf{se marchará} y el camarero pasa al siguiente cliente. \\
    \hline
  \end{tabularx}
\end{table}

\subsection{Historias de usuario}
\begin{table}[H]
  \centering
  \begin{tabularx}{\linewidth}{|l|Y|Y|Y|}
    \hline
    \textbf{ID} & \textbf{Como...} & \textbf{Quiero...} & \textbf{Para...} \\
    \hline
    HU-01 & Cliente & Pedir mi café y esperar un tiempo límite & No perder el tiempo si no me atienden me marcho. \\
    \hline
    HU-02 & Camarero & Atender al cliente que se le asigno & Atender por orden de llegada y entregar su café. \\
    \hline
    HU-03 & Barista & Preparar los cafés si no hay ninguno en la lista & Cuando no haya cafés, los hará y notificara al camarero para para que los pueda entregar si hay clientes. \\
    \hline
    HU-04 & Usuario & Poder ver el estado del programa & Saber que clientes están atendidos y por quien y saber quien recibe su café o se marcha. \\
    \hline
  \end{tabularx}
\end{table}
\clearpage

\section{Arquitectura del sistema}
\subsection{Arquitectura}
\noindent La arquitectura de está aplicación se explicará en base a lo realizado en \texttt{JavaFX} ya que sigue una arquitectura \textbf{Modelo–Vista–Controlador (MVC)} sencillo que corre únicamente en local.
\begin{itemize}
    \item \textbf{Modelo:} cuenta con las clases \texttt{cliente}, \texttt{camarero}, \texttt{barista}, \texttt{buffer}  con la lógica, el barista prepara los cafés, los camareros los reparten cuando están listos al cliente que están atendiendo y notifican cuando terminan y el cliente espera su café hasta que se lo sirvan o se agote su tiempo de espera y se marche.
    \item \textbf{Vista}: muestra los \texttt{clientes} y \texttt{camareros} con sus estados si están esperando, atendidos, si les sirvieron o si se fueron y si están libres o atendiendo un cliente y también los \textbf{cafés preparados} por el barista que están el la lista del buffer a la espera de que llegue un cliente y el camarero se lo entregué.
    \item \textbf{Controlador:} Responde al \texttt{botón de iniciar la simulación} y \texttt{actualiza los contenedores} para mostrar el estado de los clientes, camareros y cafés preparados en la interfaz.
\end{itemize}

\begin{figure}[H]
    \centering
    \includegraphics[width=1\linewidth]{MVC.png}
    \caption{Arquitectura MVC}
    \label{fig:arquitecture}
\end{figure}

\clearpage{}

\subsection{Clases}
\noindent En este caso tenemos 4 clases principales la de \textbf{Clientes}, \textbf{Camareros}, \textbf{Barista} y la de \textbf{Buffer}:

\subsection*{Cliente/Client}
\noindent \textbf{Client} representa a un cliente que entra en la cafetería, pide su café y espera un tiempo máximo (\texttt{waitingTime}) antes de marcharse si no lo recibe.

\paragraph{Atributos:}
\begin{itemize}
  \item \textbf{Identificación:} \texttt{name} (nombre del cliente), \texttt{waitingTime} (tiempo que está dispuesto a esperar en ms).
  \item \textbf{Estados:} \texttt{inFile} (si está en la cola), \texttt{assigned} (si está atendido), \texttt{served} (si ya le sirvieron el café), \texttt{left} (si ya se marchó).
  \item \textbf{Tiempos:} \texttt{arriveAt} (instante en el que llega), \texttt{timeStart} (inicio del tiempo de espera), \texttt{preparationTime} (cuanto tardo el camarero en hacer su café).
\end{itemize}

\paragraph{Flujo:}
Al inicio, el cliente marca su llegada (\texttt{timeStart}, \texttt{arriveAt}, \texttt{inFile = true}) y entra en un bucle esperando a que \texttt{served} pase a \texttt{true}. Mientras tanto, se comprueba si el tiempo de espera actual supera a \texttt{waitingTime}. Si esto pasa:
\begin{itemize}
  \item La variable para indicar que se marcho pasa a \texttt{left = true}.
  \item El hilo del cliente termina, ya que se marcho de la cafetería.
\end{itemize}

\noindent Si el camarero le entregó el café a tiempo (\texttt{served = true} antes de que se agote \texttt{waitingTime}), el cliente termina y muestra el mensaje con el \texttt{preparationTime}.

\vspace{1em}
\subsection*{Camarero/Waiter}
\noindent \textbf{Waiter} representa a un camarero que atiende solamente a un cliente asignado cada vez, entregandole el café al cliente cuando esté ya esté preparado.

\paragraph{Atributos:}
\begin{itemize}
  \item \textbf{Identificación:} \texttt{name} (nombre del camarero).
  \item \textbf{Control:} \texttt{assignament} (cliente asignado actual), \texttt{running} (permite poder parar el bucle).
  \item \textbf{Buffer:} \texttt{buffer} para que el camarero tome los cafés preparados.
  \item \textbf{Temporización:} \texttt{timeStart} (para medir la duración desde que empieza con la preparación).
\end{itemize}

\paragraph{Flujo:}
Al inicio, el camarero entra en un bucle continuo mientras \texttt{running} sea \texttt{true} o tenga un \texttt{assignament} pendiente:
\begin{itemize}
  \item Si no tiene un cliente asignado (\texttt{assignament == null}):
  \begin{itemize}
    \item Sigue esperando un poco (\texttt{Thread.sleep(15)}).
    \item Comprueba si el \texttt{CoffeController} ya le ha asignado un cliente.
  \end{itemize}
  \item Si tiene un cliente asignado que se ha marchado (\texttt{client.left}):
  \begin{itemize}
    \item Se libera la asignación (\texttt{assignament = null}) y queda disponible para atender a otro cliente.
  \end{itemize}
  \item Si tiene un cliente asignado que sigue esperando:
  \begin{itemize}
    \item Registra el \texttt{timeStart}.
    \item Obtiene un café del \texttt{Buffer} mediante \texttt{buffer.takeCoffe()}.
    \item Calcula el \texttt{preparationTime} para saber cuanto tardo en preparar/entrgar el café.
  \end{itemize}
\end{itemize}

\noindent Tras obtener el café preparado por el barista, hará lo siguiente:
\begin{itemize}
  \item Si el cliente no se marcho:
  \begin{itemize}
    \item La variable para indicar que está servido pasa a \texttt{client.served = true}.
    \item Y le pasa el \texttt{preparationTime} al cliente para mostrar cuánto tardó en recibir su café.
  \end{itemize}
  \item Si el cliente se marcho mientras esperaba el café:
  \begin{itemize}
    \item Se muestra un mensaje que informa que el cliente se marchó sin su café.
  \end{itemize}
  \item Al finalizar:
  \begin{itemize}
    \item El camarero se liberá \texttt{assignament = null} para poder atender a otro cliente.
  \end{itemize}
\end{itemize}

\noindent El método \texttt{setBuffer(Buffer buffer)} nos permite actualizar el \texttt{Buffer} por si agregamos un nuevo camarero para que el barista haga un café más.

\vspace{1em}
\subsection*{Barista}
\noindent \textbf{Barista} representa al trabajador encargado de preparar los cafés de forma continua.

\paragraph{Atributos:}
\begin{itemize}
  \item \textbf{Buffer:} \texttt{buffer} (referencia al objeto \texttt{Buffer} donde se almacenan los cafés preparados).
  \item \textbf{Control de ejecución:} \texttt{running} indica si el barista debe seguir trabajando o puede detenerse.
  \item \textbf{Simulación:} \texttt{random} (instancia de \texttt{Random}) para generar tiempos de preparación variables y hacer la simulación más realista.
\end{itemize}

\paragraph{Flujo:}
Al inicio, el barista entra en un bucle continuo mientras \texttt{running} sea \texttt{true}:
\begin{itemize}
  \item Simula el tiempo de preparación del café con \texttt{Thread.sleep()} en el método \texttt{prepareCoffe()}.
  \item Añade el café que preparo al buffer mediante \texttt{buffer.addCoffe(coffe)}.
  \item Si el \texttt{Buffer} está lleno, el barista se bloquea hasta que se liberé.
\end{itemize}

\vspace{1em}
\subsection*{Buffer}
\noindent \textbf{Buffer} es donde se guardan los cafés ya preparados por el barista y que los camareros recogerán.

\paragraph{Atributos:}
\begin{itemize}
  \item \textbf{Cola de cafés:} \texttt{coffesList} lista que almacena los cafés preparados.
  \item \textbf{Capacidad:} \texttt{capacity} indica el número máximo de cafés que puede haber preparados a la vez. Se calcula segun el número de camareros que haya.
\end{itemize}

\paragraph{Métodos principales:}
\begin{itemize}
  \item \texttt{setCapacity(int newCapacity)}: para poder actualizar la capacidad del buffer. Llama a \texttt{notifyAll()} para avisar a los hilos que puedan estar esperando por cafés.
  \item \texttt{addCoffe(String coffe)}: método que usa el \textbf{Barista} para añadir un café al buffer.
  \begin{itemize}
      \item Si la cola está llena (\texttt{coffesList.size() == capacity}), bloqueamos al barista con \texttt{wait()} hasta que un camarero consuma un café.
      \item Cuando añade el café, se llama a \texttt{notifyAll()} para avisar a los camareros que estén esperando cafés.
  \end{itemize}
  \item \texttt{takeCoffe()}: método que usan los \textbf{Camareros} para recoger un café.
  \begin{itemize}
      \item Si la cola está vacía, bloqueamos al camarero con \texttt{wait()} hasta que el barista prepare un nuevo café.
      \item Cuando hay cafés disponibles, se sacamos el primero de la lista.
      \item Tras recoger un café se llama a \texttt{notifyAll()} para notificar al barista de que vuelve a haber espacio y que preparé un café.
  \end{itemize}
  \item \texttt{getCoffesList()}: devuelve una copia de la lista de cafés actuales en el buffer (para mostrar estado de la lista de los cafés en la interfaz).
\end{itemize}

\clearpage{}

\subsection{Patrones de diseño}
\noindent Ahora vamos a explicar el \textbf{patrón de diseño empleado}, que nos permiten estructurar el flujo de ejecución y como interactan.
Usar los patrones nos ayuda a mantener el \textbf{código modular}, \textbf{legible} y \textbf{ampliarlo facilmente} en un futuro.

\subsubsection{- Patrón Productor-Consumidor}
\begin{itemize}
    \item \textbf{Productor}: genera elementos y los coloca en un \textit{buffer} compartido.
    \item \textbf{Consumidores}: sacan los elementos del buffer cuando lo necesiten.
    \item \textbf{Buffer}: lista que almacena los elementos.
\end{itemize}

En nuestra aplicación lo aplicamos de la siguiente manera:

\begin{itemize}
    \item El \textbf{Barista} es el \textbf{productor}, que prepara los cafés
    y los añade al \texttt{Buffer} cuando tenga espacio disponible.
    \item Los \textbf{Camareros} son los \textbf{consumidores}, sacando los cafés
    del \texttt{Buffer} para entregarlos a los clientes que están atendiendo.
    \item El \textbf{Buffer} es la lista compartida con los cafés preparados.
\end{itemize}

La sincronización entre hilos se realiza con los métodos \texttt{wait()} y \texttt{notifyAll()}:

\begin{itemize}
    \item \textbf{Buffer completo}, el \textbf{Barista} se bloquea hasta que algún camarero
    saque cafés de la lista.
    \item \textbf{Buffer vacío}, los \textbf{Camareros} se bloquean hasta que el barista
    prepare un nuevo café.
\end{itemize}

Esto nos garantiza:
\begin{itemize}
    \item No se preparen más cafés de los que los camareros pueden gestionar (control de capacidad).
    \item Los camareros no intenten entregar un café que todavía no existe.
\end{itemize}

\begin{figure}[H]
    \centering
    \includegraphics[width=.85\linewidth]{productor-consumidor.png}
    \caption{Patrón productor/consumidor}
    \label{fig:patterns}
\end{figure}

\subsection{Stack empleado}
\noindent Para el desarrollo nuestra aplicación tanto en terminal como con interfaz, se utilizaron las siguietes tecnológias:
\begin{itemize}
    \item Utilizamos \texttt{Java} como \textbf{entorno de ejecución} principal.
    \item Utilizamos \texttt{JavaFX} para crear una \textbf{interfaz con Java}.
\end{itemize}

\subsection{Interfaz}
\noindent Se creo una interfaz muy simple en \texttt{JavaFX}, para \textbf{mostrar el programa de forma clara y bonita}, para saber cuando son atendidos los clientes y cuanto tardan en recibir su café o si se han ido. Como podemos apreciar tenemos un \textbf{título} y debajo el \textbf{botón} para comenzar con esta simulación del servicio de la cafetería y luego tenemos \textbf{tres contenedores}
uno para clientes, uno para los camareros donde mostramos su estado y un contenedor para los cafés preparados por el barista que están en cola de espera.

\begin{figure}[H]
    \centering
    \includegraphics[width=1\linewidth]{interfaz.png}
    \caption{Interfaz en JavaFX}
    \label{fig:interfaz}
\end{figure}

\end{document}